% !TEX root = ../main.tex

\chapter{Validazione dell’artefatto}
\label{chap:evaluation}

Questo capitolo presenta la validazione dell’artefatto PENTION-M secondo l’approccio di Design Science Research adottato.
La validazione non è intesa come un singolo test finale, ma come un processo integrato e progressivo (\emph{validation by construction}), che combina valutazione tecnica, riscontri qualitativi da parte degli stakeholder e un \emph{experience report} sulle iterazioni di sviluppo.

\section{Inquadramento della validazione secondo il Design Science Research}
\label{sec:validation_dsr}

Nel paradigma della \emph{Design Science Research} (DSR), la validazione dell’artefatto
non è concepita come una fase finale isolata, ma come una componente intrinseca
dell’intero processo di progettazione e sviluppo.
A differenza di approcci sperimentali tradizionali, in cui la valutazione è spesso
limitata alla misurazione delle prestazioni di singoli modelli o algoritmi,
il DSR pone l’accento sulla capacità dell’artefatto di rispondere in modo coerente
e verificabile a un problema reale, tenendo conto del contesto applicativo,
dei vincoli operativi e delle esigenze degli stakeholder.

Nel caso di \textsc{PENTION-M}, la validazione è stata affrontata secondo un approccio
di \emph{validation by construction}, già delineato nel Capitolo~\ref{chap:dsr}.
Tale approccio implica che le proprietà di qualità dell’artefatto
— robustezza, coerenza fisica, auditabilità e usabilità operativa —
non vengano valutate esclusivamente attraverso un esperimento finale,
ma emergano progressivamente dalle scelte metodologiche,
dalle iterazioni progettuali e dalla struttura stessa dell’architettura implementata.

\subsection*{Validazione dell’artefatto vs valutazione di singoli modelli}

Un elemento centrale per l’inquadramento della validazione riguarda la distinzione
tra la valutazione di singoli modelli e la validazione dell’artefatto nel suo complesso.
Nel contesto di questa tesi, l’artefatto non coincide con un singolo modello di
Machine Learning, bensì con un sistema cyber--fisico simulato che integra
modellazione fisica, modelli di \emph{Physics-Informed Machine Learning},
pipeline MLOps e meccanismi di logging forense in un flusso end-to-end coerente.

La valutazione di un modello isolato (ad esempio in termini di accuratezza,
RMSE o F1-score) fornisce informazioni utili ma parziali,
insufficienti a caratterizzare il comportamento del sistema in uno scenario
operativo reale.
Al contrario, la validazione dell’artefatto richiede di verificare proprietà
sistemiche, quali:
(i) la coerenza tra input fisici, inferenze e output prodotti;
(ii) la stabilità del comportamento a runtime;
(iii) la tracciabilità e riproducibilità delle decisioni;
(iv) l’aderenza ai workflow investigativi e ai requisiti forensi.

In \textsc{PENTION-M}, tali proprietà non sono demandate a un singolo componente,
ma distribuite lungo l’architettura a layer descritta nel
Capitolo~\ref{chap:architecture}.
La validazione assume quindi una dimensione trasversale,
che attraversa l’intera pipeline e coinvolge simultaneamente
scelte modellistiche, architetturali e operative.

\subsection*{Concetto di \emph{validation by construction} in sistemi cyber--fisici}

Il concetto di \emph{validation by construction} risulta particolarmente appropriato
nel contesto dei sistemi cyber--fisici e, in particolare, di quelli destinati
a impieghi forensi e investigativi.
In tali sistemi, la correttezza non può essere valutata unicamente sulla base
dell’output finale, ma deve essere garantita anche rispetto al processo che conduce
alla decisione.

Nel progetto \textsc{PENTION-M}, la validazione by construction si manifesta attraverso:
\begin{itemize}
  \item l’integrazione esplicita di conoscenza fisica nei modelli di inferenza,
  al fine di limitare soluzioni numericamente corrette ma fisicamente implausibili;
  \item la separazione architetturale delle responsabilità tra simulazione,
  inferenza, monitoring e logging forense;
  \item la produzione nativa di metadati, metriche runtime ed evidenze verificabili
  come parte integrante dell’esecuzione della pipeline;
  \item l’allineamento continuo tra requisiti emersi dalla campagna di
  Requirements Engineering e le iterazioni dell’artefatto.
\end{itemize}

In questo quadro, la validazione non è un’attività aggiuntiva o retrospettiva,
ma una proprietà emergente del modo in cui il sistema è stato progettato e realizzato.
Le scelte architetturali illustrate nel Capitolo~\ref{chap:architecture}
— in particolare l’introduzione di layer dedicati al monitoring,
al ModelOps e al logging forense —
rappresentano esempi concreti di come i criteri di validazione siano stati
incorporati direttamente nella struttura dell’artefatto.

\subsection*{Collegamento tra requisiti, architettura e criteri di validazione}

Un aspetto distintivo dell’approccio adottato è il collegamento esplicito
tra requisiti, architettura e criteri di validazione.
I requisiti emersi dalla campagna di Requirements Engineering
(Sezione~\ref{sec:requirements_engineering})
non sono stati trattati come semplici vincoli iniziali,
ma come riferimenti costanti per guidare e valutare l’evoluzione dell’artefatto.

Ogni iterazione di \textsc{PENTION-M} ha tradotto specifici requisiti
— operativi, tecnici, fisici e forensi —
in scelte progettuali verificabili,
come descritto nella Sezione~\ref{sec:artifact_iterations}.
Di conseguenza, la validazione dell’artefatto consiste nel verificare
la coerenza tra:
(i) i requisiti dichiarati;
(ii) le soluzioni architetturali adottate;
(iii) il comportamento osservato del sistema durante l’esecuzione end-to-end.

Il presente capitolo si colloca quindi come naturale prosecuzione
del percorso metodologico e architetturale già delineato.
Le sezioni successive presenteranno la validazione dell’artefatto
attraverso il confronto con gli stakeholder e l’analisi tecnica
del comportamento del sistema,
mostrando come le proprietà dichiarate emergano in modo coerente
dal processo di progettazione seguito.

\section{Validazione tramite stakeholder}
\label{sec:validation_stakeholders}

Nel contesto della \emph{Design Science Research}, la validazione tramite stakeholder
non coincide con una nuova fase di elicitazione dei requisiti,
già descritta nel Capitolo~\ref{chap:dsr},
ma rappresenta una forma di \emph{valutazione ex post} dell’artefatto sviluppato.
In questa prospettiva, gli stakeholder non sono considerati come fonte di input progettuali,
bensì come riferimento esterno per verificare la coerenza, la credibilità e
la rilevanza operativa del sistema risultante.

Per il progetto \textsc{PENTION-M}, il confronto con stakeholder operativi,
industriali e accademici è stato utilizzato come strumento di validazione qualitativa
dell’artefatto nel suo complesso.
L’obiettivo non è stato quello di misurare prestazioni puntuali,
ma di verificare se le scelte architetturali e metodologiche adottate
risultassero plausibili, utili e giustificabili rispetto al contesto
forense-investigativo di riferimento.

Questa forma di validazione è coerente con l’approccio di
\emph{validation by construction} discusso nella Sezione~\ref{sec:validation_dsr}:
le proprietà del sistema vengono valutate non come risultati isolati,
ma come conseguenza diretta del modo in cui l’artefatto è stato progettato,
implementato e integrato.

\subsection{Feedback delle Law Enforcement Agencies}
\label{subsec:validation_lea}

Il primo livello di validazione tramite stakeholder riguarda il confronto
con le \emph{Law Enforcement Agencies} (LEA),
che rappresentano il principale contesto applicativo di riferimento per PENTION-M.
In questa fase, l’artefatto sviluppato è stato valutato in termini di
coerenza con i workflow investigativi reali,
utilità percepita delle funzionalità offerte
e adeguatezza delle modalità di presentazione degli output.

Dal punto di vista operativo, PENTION-M risulta coerente con le modalità
di intervento tipiche delle LEA,
in particolare negli scenari di pattugliamento urbano,
controllo di aree logistiche e ispezione di ambienti chiusi.
La pipeline end-to-end, descritta nel Capitolo~\ref{chap:architecture},
supporta un flusso di lavoro in cui l’analisi automatizzata
fornisce indicazioni rapide e interpretabili,
senza richiedere competenze specialistiche in chimica o modellazione fisica.

Un aspetto rilevante emerso dal confronto con il dominio LEA riguarda
l’utilità percepita della localizzazione della sorgente e della classificazione NPS
come strumenti di supporto decisionale.
In particolare, la combinazione tra:
(i) indicazione di presenza/assenza del segnale,
(ii) stima spaziale della possibile origine,
e (iii) livello di confidenza associato all’output,
risulta allineata con le pratiche operative attuali,
in cui l’AI è utilizzata come ausilio all’interpretazione
e non come sostituto della decisione umana.

Infine, un elemento di forte rilevanza per la validazione da parte delle LEA
è rappresentato dalla tracciabilità forense del processo decisionale.
La possibilità di ricostruire a posteriori
input, parametri di contesto, versioni dei modelli e output prodotti
contribuisce a rendere il sistema compatibile
con esigenze di auditabilità, catena di custodia e accountability,
aspetti fondamentali in ambito investigativo.

\subsection{Feedback dei partner industriali e di progetto}
\label{subsec:validation_industry}

Il secondo livello di validazione tramite stakeholder riguarda
il confronto con partner industriali e di progetto,
il cui interesse principale è rivolto alla robustezza architetturale,
alla modularità del sistema e al potenziale di integrazione futura.

Dal punto di vista architetturale, PENTION-M è stato valutato positivamente
per la separazione esplicita delle responsabilità tra i diversi componenti:
ingestion, simulazione fisica, modelli PIML, classificazione,
monitoring e logging forense.
Questa organizzazione a microservizi containerizzati,
descritta nel Capitolo~\ref{chap:architecture},
favorisce sia la manutenibilità del sistema
sia l’adattabilità a contesti operativi differenti.

Un ulteriore elemento di validazione riguarda la neutralità tecnologica
dell’architettura.
L’adozione di interfacce ben definite tra i moduli consente,
almeno a livello concettuale,
la sostituzione o l’estensione di singoli componenti
(ad esempio modelli di classificazione o servizi di simulazione)
senza richiedere una riprogettazione complessiva della pipeline.
Questa caratteristica è particolarmente rilevante
in un contesto di progetto europeo,
in cui interoperabilità e riusabilità rappresentano requisiti chiave.

Nel complesso, il feedback dei partner industriali
conferma che PENTION-M non è un prototipo monolitico,
ma un’architettura coerente e scalabile,
potenzialmente integrabile in ecosistemi più ampi
di sistemi di analisi e supporto investigativo.

\subsection{Feedback accademico e supervisione scientifica}
\label{subsec:validation_academic}

Il terzo livello di validazione tramite stakeholder
è rappresentato dalla supervisione accademica e dal confronto scientifico,
che ha avuto come obiettivo principale
la verifica della solidità metodologica e concettuale dell’approccio adottato.

Dal punto di vista scientifico,
l’integrazione del \emph{Physics-Informed Machine Learning}
all’interno di una pipeline MLOps simulata
è stata valutata come coerente con lo stato dell’arte
e adeguata al dominio applicativo considerato.
In particolare, l’uso di vincoli fisici e modelli di dispersione
come guida per l’inferenza contribuisce a limitare
comportamenti non plausibili dei modelli data-driven,
rafforzando l’interpretabilità e la robustezza del sistema.

La scelta metodologica di adottare il paradigma
di \emph{Design Science Research} è stata ritenuta appropriata
per un progetto che combina ricerca applicata,
sviluppo ingegneristico e confronto con stakeholder reali.
La distinzione tra requisiti, progettazione dell’artefatto
e validazione ex post consente di mantenere
una struttura rigorosa e difendibile del lavoro di tesi.

Infine, particolare attenzione è stata posta
agli aspetti di auditabilità, riproducibilità e trasparenza.
La presenza di meccanismi di logging forense,
versionamento dei modelli e tracciabilità end-to-end
rappresenta non solo un requisito operativo,
ma anche un contributo metodologico rilevante
per l’adozione dell’AI in contesti forensi e investigativi.

Nel loro insieme, questi elementi confermano che PENTION-M
può essere considerato un artefatto valido
dal punto di vista scientifico, metodologico e applicativo,
pur operando in un contesto simulato.

\section{Validazione tecnica dell’artefatto}
\label{sec:technical_validation}

La validazione tecnica di \textsc{PENTION-M} ha lo scopo di verificare
il comportamento dell’artefatto nel suo insieme,
considerato come sistema cyber-fisico end-to-end,
piuttosto che come collezione di modelli valutati in modo isolato.
In linea con il paradigma di \emph{Design Science Research}
e con l’approccio di \emph{validation by construction},
le attività di validazione sono state progettate per
riflettere condizioni operative realistiche,
pur operando in un contesto interamente simulato.

Le analisi tecniche presentate in questa sezione
sono state condotte utilizzando una specifica campagna di validazione,
raccolta nella cartella \texttt{validation/} del progetto.
Tale campagna include test end-to-end automatizzati,
stress test della pipeline MLOps
e verifiche di coerenza degli output prodotti dal sistema
durante l’esecuzione completa della pipeline.

È importante sottolineare che questa fase di validazione
non mira a dimostrare l’ottimalità assoluta dei singoli modelli,
bensì a verificare che l’architettura progettata
soddisfi requisiti fondamentali di:
(i) eseguibilità end-to-end,
(ii) stabilità operativa,
(iii) rispetto dei vincoli temporali,
(iv) coerenza dei risultati intermedi,
(v) tracciabilità delle esecuzioni.
La valutazione delle prestazioni specifiche dei modelli PIML,
della localizzazione della sorgente e della classificazione NPS
è trattata separatamente nelle sezioni successive.

\subsection{Validazione delle prestazioni runtime}
\label{subsec:runtime_validation}

La validazione delle prestazioni runtime si concentra
sull’analisi del comportamento temporale della pipeline end-to-end,
con particolare attenzione alla latenza complessiva
e alla compatibilità del sistema con uno scenario operativo mobile.
In un contesto investigativo sul campo,
la capacità di fornire output in tempi contenuti
rappresenta un requisito essenziale,
indipendentemente dalla qualità dei modelli sottostanti.

La pipeline di \textsc{PENTION-M} è stata validata
mediante test end-to-end automatizzati,
in cui un payload simulato completo
(vettori spettrali EI, parametri ambientali,
coordinate GPS del sensore e sorgente)
viene inviato al servizio di ingestion
e propagato attraverso tutti i layer del sistema.
Durante questa esecuzione, il sistema misura e registra
la latenza complessiva e le metriche di monitoraggio associate.

I test end-to-end verificano innanzitutto
la corretta esecuzione della pipeline,
controllando che:
\begin{itemize}
  \item l’ingestion dei dati venga completata senza errori;
  \item i servizi di simulazione fisica, correzione PIML,
        localizzazione e classificazione producano output validi;
  \item il bundle forense venga generato correttamente;
  \item le metriche di monitoring siano disponibili a valle dell’esecuzione.
\end{itemize}

Dal punto di vista temporale,
la latenza end-to-end viene misurata come intervallo
tra l’invio del payload al servizio di ingestion
e la ricezione della risposta finale,
che include sia i risultati di inferenza
sia il blocco di monitoring.
Nei test automatizzati,
la pipeline è considerata conforme
quando la latenza complessiva rimane al di sotto
di una soglia compatibile con l’uso operativo,
definita a priori durante la progettazione del sistema.

Oltre ai test funzionali singoli,
la validazione runtime include una campagna di stress test,
in cui la pipeline viene eseguita ripetutamente
con input variabili (posizione, meteo, spettri rumorosi),
simulando una sequenza di eventi operativi consecutivi.
Questi test consentono di osservare:
(i) la stabilità delle latenze,
(ii) l’assenza di degrado progressivo delle prestazioni,
(iii) il comportamento dei servizi containerizzati
sotto carico ripetuto.

Durante gli stress test,
vengono inoltre raccolte metriche di utilizzo delle risorse
dei principali container della pipeline
(ingestion, monitoring, simulazione fisica,
correzione PIML, localizzazione e classificazione),
al fine di verificare che il costo computazionale
rimanga compatibile con un’architettura mobile
basata su hardware limitato.
Questa analisi non ha l’obiettivo di ottimizzare le prestazioni,
ma di escludere configurazioni intrinsecamente non sostenibili
in uno scenario di laboratorio mobile su veicolo.

Nel complesso, i risultati dei test end-to-end e degli stress test
confermano che la pipeline di \textsc{PENTION-M}
è in grado di operare in modo stabile e ripetibile,
fornendo output completi e monitorabili
entro tempi compatibili con un contesto operativo simulato.
Questa evidenza costituisce un prerequisito fondamentale
per le successive validazioni di coerenza fisica,
accuratezza spaziale e affidabilità della classificazione,
che vengono analizzate nelle sezioni seguenti.

\subsection{Validazione dei modelli PIML}
\label{sec:val_piml}

Questo sottocapitolo presenta la validazione tecnica del servizio di \emph{correzione PIML della dispersione} (CorrectionDispersion--PIML), con l’obiettivo di verificare che le mappe prodotte dal modello risultino (i) coerenti rispetto a una \emph{ground truth} simulata di riferimento e (ii) compatibili con vincoli fisici e geometrici dello scenario urbano (in particolare la presenza di aree edificate).

In coerenza con l’impostazione \emph{by construction} della tesi, la validazione qui descritta non mira a “massimizzare una metrica” in astratto, ma a controllare proprietà operative e fisiche che rendono l’output difendibile e interpretabile nel contesto architetturale di \textsc{PENTION-M}.

\subsubsection{Setup sperimentale e pipeline di validazione}
\label{sec:val_piml_setup}

La validazione è stata eseguita tramite lo script/notebook \texttt{validation/PIML/PIML\_Validation\_All.ipynb}, che implementa un ciclo di test ripetuto su un insieme di scenari simulati.
Per ogni simulazione:
\begin{itemize}
    \item viene caricata una mappa di riferimento della dispersione ``reale'' (ground truth) dal repository dei dati simulati (\texttt{real\_dispersion});
    \item viene generata o recuperata la mappa di dispersione di baseline (gaussiana) e inviata al servizio di correzione PIML tramite una chiamata API (funzione helper in \texttt{validation/PIML/utils\_piml.py});
    \item si ottiene in output una mappa di concentrazione corretta, della stessa dimensionalità della mappa di input;
    \item si calcola un set di metriche fisicamente motivate (definite in \texttt{validation/PIML/metrics.py});
    \item i risultati vengono aggregati e salvati in un file strutturato (\texttt{validation/PIML/validation\_results\_piml.csv}).
\end{itemize}

L’obiettivo del processo è quindi valutare la qualità delle mappe corrette rispetto alla ground truth e rispetto a vincoli geometrici e direzionali (vento + maschera edifici), \emph{senza} introdurre assunzioni ulteriori non presenti nella pipeline di progetto.

\subsubsection{Metriche di validazione}
\label{sec:val_piml_metrics}

La valutazione utilizza quattro metriche, scelte per catturare aspetti complementari dell’output:

\paragraph{Errore RMSE su area libera (\texttt{rmse\_free}).}
La metrica misura l’errore quadratico medio tra mappa predetta e mappa di riferimento, calcolato esclusivamente sulle celle non edificate (\emph{free space}), utilizzando la maschera urbana.
Formalmente, indicando con $M$ la maschera binaria degli edifici e con $\Omega=\{(i,j): M_{ij}=0\}$ l’insieme delle celle libere:
\[
\mathrm{RMSE}_{\mathrm{free}} =
\sqrt{\frac{1}{|\Omega|}\sum_{(i,j)\in\Omega}\left(\hat{C}_{ij}-C_{ij}\right)^2}.
\]
Questa scelta evita che l’errore sia dominato da regioni in cui (per costruzione) la concentrazione dovrebbe risultare nulla o fortemente penalizzata.

\paragraph{Indice di smoothness (\texttt{smoothness}).}
La smoothness è un indicatore di regolarità spaziale della mappa, costruito a partire dai gradienti discreti; l’obiettivo è evidenziare la presenza di artefatti ad alta frequenza (es. discontinuità non plausibili) e caratterizzare la “stabilità spaziale” del campo predetto.
In pratica, la metrica implementa un proxy di \emph{total variation} basato sulle derivate discrete.

\paragraph{Allineamento al vento (\texttt{wind\_alignment}).}
Questa metrica valuta la coerenza direzionale del campo rispetto al vento dominante.
A partire dal gradiente $\nabla \hat{C}$ e dal vettore del vento $\mathbf{v}=[\cos(\theta),\sin(\theta)]$, viene calcolato un punteggio medio normalizzato basato sul prodotto scalare tra gradiente e direzione del vento (con normalizzazione per evitare dipendenze dalla scala del gradiente):
\[
s = \mathbb{E}\left[\frac{\nabla \hat{C}\cdot \mathbf{v}}{\|\nabla \hat{C}\|\cdot\|\mathbf{v}\|+\varepsilon}\right],
\qquad s\in[-1,1].
\]
Valori prossimi a $1$ indicano una maggiore coerenza direzionale tra variazione spaziale del campo e direzione imposta dal vento, mentre valori bassi (o negativi) suggeriscono campi meno compatibili con la direzionalità attesa.

\paragraph{Tasso di violazione della maschera edifici (\texttt{building\_violation}).}
La metrica stima quanta ``massa'' di concentrazione (in senso relativo) ricade in area edificata.
È implementata come rapporto tra concentrazione cumulata sulle celle edificate e concentrazione totale della mappa:
\[
\mathrm{BV}=\frac{\sum_{i,j}\hat{C}_{ij}\,M_{ij}}{\sum_{i,j}\hat{C}_{ij}+\varepsilon}.
\]
Un valore vicino a zero indica un’uscita compatibile con il vincolo geometrico (assenza di concentrazione nelle regioni edificate).

\subsubsection{Risultati aggregati su campagna di simulazioni}
\label{sec:val_piml_results}

I risultati della campagna sono salvati nel file \texttt{validation\_results\_piml.csv}, che contiene un record per simulazione con le quattro metriche descritte.
La Tabella~\ref{tab:piml_validation_summary} riporta le statistiche principali (media, deviazione standard, minimo e massimo) calcolate sull’intero set di simulazioni.

\begin{table}[H]
\centering
\renewcommand{\arraystretch}{1.15}
\begin{tabular}{lcccc}
\toprule
\textbf{Metrica} & \textbf{Media} & \textbf{Std} & \textbf{Min} & \textbf{Max} \\
\midrule
\texttt{rmse\_free}        & 0.0575 & 0.0318 & 0.0067 & 0.2079 \\
\texttt{smoothness}        & 0.0075 & 0.0030 & 0.0006 & 0.0183 \\
\texttt{wind\_alignment}   & 0.4287 & 0.1467 & 0.1132 & 0.7759 \\
\texttt{building\_violation} & 0.0000 & 0.0003 & 0.0000 & 0.0031 \\
\bottomrule
\end{tabular}
\caption{Sintesi delle metriche di validazione PIML calcolate sulle simulazioni della campagna (\texttt{validation/PIML/validation\_results\_piml.csv}).}
\label{tab:piml_validation_summary}
\end{table}

Nel complesso, la metrica \texttt{building\_violation} risulta prossima a zero nella quasi totalità delle simulazioni, evidenziando un comportamento coerente con il vincolo geometrico imposto dalla maschera urbana.
In parallelo, il valore medio di \texttt{wind\_alignment} suggerisce una coerenza direzionale apprezzabile rispetto al vento, mentre la \texttt{smoothness} rimane contenuta, indicando assenza di oscillazioni spaziali marcate nella maggior parte dei casi.
L’errore \texttt{rmse\_free} fornisce invece una misura diretta di scostamento dalla ground truth sulle aree osservabili non edificate, permettendo di caratterizzare in modo quantitativo la qualità del campo corretto.

\subsubsection{Discussione e limiti}
\label{sec:val_piml_discussion}

È importante sottolineare che la campagna di validazione è condotta in un contesto simulato controllato: la ground truth è generata coerentemente con lo scenario e le maschere urbane disponibili.
Questo garantisce riproducibilità e consente di verificare proprietà fisiche e geometriche in modo sistematico, ma non equivale a una validazione su dati reali.
Inoltre, le metriche utilizzate sono pensate per controllare \emph{plausibilità fisica} e \emph{compatibilità geometrica} dell’output, più che per ottimizzare un singolo indicatore globale.

Nel quadro di \textsc{PENTION-M}, questa validazione supporta il ruolo del CorrectionDispersion--PIML come servizio del Layer~3: l’output prodotto risulta misurabile, tracciabile e verificabile rispetto a criteri fisicamente interpretabili, costituendo una base solida per le fasi downstream della pipeline (localizzazione e reporting operativo/forense), che vengono analizzate nei sottocapitoli successivi.

\subsection{Validazione della localizzazione della sorgente}
\label{subsec:source_localization_validation}

La localizzazione della sorgente di emissione rappresenta una componente critica dell’artefatto \textsc{PENTION-M}, in quanto costituisce il collegamento diretto tra la modellazione della dispersione e il supporto operativo alle attività investigative.
Nel contesto del paradigma di \emph{Design Science Research}, la validazione di tale componente non è intesa come una valutazione isolata di un singolo modello predittivo, ma come una verifica della sua adeguatezza funzionale all’interno della pipeline cyber-fisica complessiva.

In \textsc{PENTION-M}, la stima della posizione della sorgente è affidata a un modello di \emph{Emission Source Localization} fisicamente informato, integrato a valle del modulo di correzione PIML della dispersione.
La validazione è stata condotta utilizzando il modulo \texttt{EmissionSourceLocalization\_PIML}, coerentemente con l’architettura descritta nel Capitolo~\ref{chap:architecture}, e si basa su scenari simulati in cui la posizione reale della sorgente è nota.

\paragraph{Criteri di valutazione.}
La qualità della localizzazione è stata valutata in termini di errore spaziale tra la posizione stimata e la posizione reale della sorgente.
In particolare, sono state adottate metriche standard e direttamente interpretabili in ambito operativo:
\begin{itemize}
    \item \textbf{Mean Absolute Error (MAE)} spaziale, espresso in metri;
    \item \textbf{Root Mean Square Error (RMSE)} spaziale, espresso in metri.
\end{itemize}
Tali metriche consentono di quantificare sia l’errore medio di localizzazione sia la presenza di deviazioni più marcate, fornendo una misura sintetica dell’affidabilità del sistema.

\paragraph{Stabilità delle stime spaziali.}
I risultati aggregati della validazione, calcolati su un insieme di scenari simulati, mostrano una buona stabilità delle stime spaziali.
L’errore medio assoluto si colloca nell’ordine di alcune decine di metri, mentre l’RMSE evidenzia una dispersione contenuta degli errori attorno al valore medio.
Questo comportamento indica che il modello di localizzazione produce stime consistenti e non affette da instabilità numeriche o da outlier sistematici.

Dal punto di vista del sistema complessivo, tale stabilità è particolarmente rilevante, poiché la localizzazione non è utilizzata come valore puntuale assoluto, ma come informazione di supporto alla ricostruzione del contesto operativo e investigativo.

\paragraph{Robustezza rispetto a rumore e osservabilità parziale.}
La validazione è stata condotta in condizioni che riflettono implicitamente la presenza di rumore e osservabilità parziale, tipiche degli scenari simulati di dispersione urbana.
La stima della sorgente avviene infatti a partire da mappe di concentrazione corrette tramite modelli PIML, che incorporano incertezza legata a condizioni ambientali, configurazione dei sensori e geometria urbana.

I risultati ottenuti mostrano che, pur in presenza di tali fattori di incertezza, il modello di localizzazione mantiene un comportamento robusto, senza degradazioni improvvise delle prestazioni.
Ciò suggerisce che l’approccio fisicamente informato contribuisce a vincolare lo spazio delle soluzioni ammissibili, riducendo l’impatto del rumore sulle stime finali.

\paragraph{Ordine di grandezza degli errori spaziali.}
L’ordine di grandezza degli errori di localizzazione osservati è coerente con le aspettative per un sistema di rilevamento mobile operante in ambiente urbano.
Errori nell’ordine delle decine di metri risultano compatibili con:
\begin{itemize}
    \item la natura indiretta del problema di localizzazione;
    \item l’uso di dati simulati affetti da rumore e incompletezza;
    \item l’obiettivo di fornire un’indicazione spaziale utile al supporto decisionale, piuttosto che una localizzazione metrica assoluta.
\end{itemize}

In un contesto investigativo reale, tali stime sono destinate a essere integrate con ulteriori informazioni (osservazioni sul campo, dati contestuali, esperienza degli operatori), rafforzando il ruolo della localizzazione come strumento di \emph{decision support}.

\paragraph{Discussione nel contesto DSR.}
Nel quadro della validazione \emph{by construction} adottata in questa tesi, i risultati della localizzazione della sorgente non vengono interpretati come una dimostrazione di ottimalità del modello, bensì come evidenza della sua adeguatezza funzionale all’interno dell’artefatto \textsc{PENTION-M}.
La coerenza delle stime, la stabilità degli errori e l’ordine di grandezza delle prestazioni contribuiscono a confermare che il modulo di localizzazione soddisfa i requisiti operativi e metodologici definiti nella fase di Requirements Engineering, completando in modo coerente il Layer di inferenza fisicamente informata della pipeline.

\subsection{Validazione della classificazione NPS}
\label{subsec:nps_validation}

La validazione del modulo di classificazione delle \emph{New Psychoactive Substances} (NPS) è stata condotta con l’obiettivo di verificare l’affidabilità operativa del classificatore impiegato in \textsc{PENTION-M}, piuttosto che di massimizzarne le prestazioni su metriche puramente accademiche.
In linea con il paradigma di \emph{validation by construction}, l’analisi si concentra sulla qualità delle confidenze prodotte, sulla loro calibrazione e sulla robustezza del modello rispetto a perturbazioni realistiche dei dati di ingresso.

Il classificatore validato è un modello \textbf{XGBoost} addestrato su spettri EI vettorizzati, utilizzato come componente principale della pipeline decisionale.
L’output del modello consiste in una predizione di classe accompagnata da una stima di confidenza, successivamente utilizzata per supportare la decisione dell’operatore umano.

\paragraph{Calibrazione delle confidenze}
Un primo aspetto critico analizzato riguarda la calibrazione delle probabilità prodotte dal modello.
In contesti forensi e investigativi, una confidenza numerica non calibrata può risultare fuorviante, inducendo un’eccessiva fiducia in predizioni errate o marginali.

La Figura~\ref{fig:nps_calibration_curve} mostra la curva di calibrazione ottenuta dopo l’applicazione di una strategia di \emph{Temperature Scaling} con parametro $T=1.8$.
La curva mette in relazione la confidenza predetta dal modello con l’accuratezza empirica osservata, consentendo di valutare lo scostamento rispetto al comportamento ideale di perfetta calibrazione.

\begin{figure}[H]
\centering
\includegraphics[width=0.6\textwidth]{calibration_curve.png}
\caption{Curva di calibrazione del classificatore XGBoost dopo Temperature Scaling ($T=1.8$).}
\label{fig:nps_calibration_curve}
\end{figure}

I risultati mostrano un buon allineamento tra confidenza predetta e accuratezza reale, in particolare nella regione di alta confidenza, che rappresenta la zona di maggiore rilevanza operativa.
Il valore di \emph{Brier score} ottenuto ($\approx 0.02$) conferma quantitativamente la buona qualità della calibrazione, indicando una ridotta discrepanza media tra probabilità predette ed esito corretto.

\paragraph{Distribuzione delle confidenze}
Per valutare il comportamento globale del classificatore, è stata analizzata la distribuzione delle confidenze associate alle predizioni.
La Figura~\ref{fig:nps_confidence_distribution} evidenzia come la maggior parte delle inferenze si collochi in una regione di confidenza elevata, con una progressiva rarefazione verso valori intermedi e bassi.

\begin{figure}[H]
\centering
\includegraphics[width=0.6\textwidth]{confidence_histogram.png}
\caption{Distribuzione delle confidenze prodotte dal classificatore XGBoost dopo Temperature Scaling.}
\label{fig:nps_confidence_distribution}
\end{figure}

Questo comportamento è coerente con uno scenario operativo in cui il sistema è chiamato a fornire indicazioni affidabili principalmente in presenza di segnali informativi sufficientemente forti, lasciando all’operatore la gestione dei casi ambigui o a bassa confidenza.

\paragraph{Robustezza rispetto al rumore spettrale}
Un ulteriore aspetto di validazione riguarda la robustezza del classificatore rispetto a perturbazioni realistiche degli spettri di massa.
A tal fine, è stato condotto un test di robustezza introducendo rumore gaussiano additivo di diversa intensità sugli spettri EI di input.

La Figura~\ref{fig:nps_noise_robustness} mostra l’andamento della confidenza media al crescere del livello di rumore.
I risultati indicano una degradazione graduale e controllata delle confidenze, senza collassi improvvisi del comportamento del modello.

\begin{figure}[H]
\centering
\includegraphics[width=0.6\textwidth]{noise_robustness.png}
\caption{Robustezza del classificatore NPS rispetto a rumore spettrale crescente.}
\label{fig:nps_noise_robustness}
\end{figure}

Anche in presenza di livelli di rumore significativi, il modello mantiene confidenze elevate e stabili, suggerendo una buona tolleranza a condizioni di acquisizione non ideali, tipiche di scenari operativi sul campo.

\paragraph{Discussione qualitativa degli errori}
L’analisi qualitativa delle predizioni errate evidenzia come gli errori residui siano principalmente associati a:
\begin{itemize}
  \item spettri fortemente degradati o incompleti;
  \item composti appartenenti a classi chimiche strutturalmente affini;
  \item segnali con intensità prossime alla soglia di rilevamento.
\end{itemize}

Tali casi rappresentano situazioni intrinsecamente ambigue anche per un analista umano, e risultano coerenti con l’impostazione del sistema come \emph{decision support tool} piuttosto che come sostituto del giudizio esperto.

Nel complesso, i risultati della validazione mostrano che il classificatore NPS integrato in \textsc{PENTION-M} fornisce predizioni stabili, confidenze calibrate e un comportamento robusto rispetto a perturbazioni realistiche, risultando adeguato per l’impiego in una pipeline investigativa mobile e forense-ready.

\subsection{Effetto del Temperature Scaling adattivo}
\label{sec:temp_scaling}

Nel contesto di sistemi decisionali a supporto investigativo, la calibrazione delle probabilità associate alle predizioni riveste un ruolo cruciale. In particolare, nel caso della classificazione di Nuove Sostanze Psicoattive (NPS), la confidenza restituita dal modello non rappresenta un mero indicatore accessorio, ma costituisce parte integrante dell’evidenza forense prodotta dal sistema. Un valore di confidenza sovrastimato può infatti indurre una falsa percezione di certezza, compromettendo la correttezza dell’interpretazione operativa e la successiva utilizzabilità legale del risultato.

Durante le prime fasi di sviluppo di \textsc{PENTION-M}, il classificatore NPS basato su XGBoost mostrava un comportamento di marcata \emph{overconfidence}, con probabilità frequentemente prossime all’unità anche in presenza di spettri EI degradati, rumorosi o acquisiti in condizioni ambientali sfavorevoli. Tale fenomeno è ampiamente documentato nella letteratura sui modelli di boosting e rappresenta una criticità ben nota nei sistemi di classificazione probabilistica non calibrati.

Per mitigare questo problema, è stata inizialmente introdotta una tecnica di \emph{Temperature Scaling} post-hoc, secondo l’approccio proposto da Guo et al.~\cite{guo2017calibration}, applicata alle probabilità prodotte dal classificatore XGBoost. In una fase preliminare, il sistema adottava un valore di temperatura fisso pari a $T = 1.8$, scelto in modo euristico al fine di ridurre i casi più evidenti di saturazione della softmax. Sebbene efficace nel contenere l’overconfidence, tale scelta non era supportata da una validazione sistematica né risultava adattabile alle diverse condizioni operative della pipeline.

Sebbene l’approccio originale preveda l’uso di un parametro di temperatura statico,
ottimizzato su un dataset di validazione, nel contesto operativo di \textsc{PENTION-M}
si è reso necessario estendere tale meccanismo verso una formulazione adattiva,
in grado di riflettere la variabilità fisica e computazionale del sistema.

Successivamente, è stata condotta una fase di studio e validazione dedicata alla calibrazione, basata su un dataset di validazione costruito ad hoc, comprendente spettri rappresentativi delle diverse classi NPS e versioni perturbate degli stessi mediante l’aggiunta di rumore controllato. Su tale dataset è stata eseguita una \emph{sweep analysis} dei valori di temperatura nell’intervallo $[1.0, 3.0]$, valutando metriche di calibrazione quali il \emph{Brier Score} e l’Expected Calibration Error (ECE). I risultati hanno evidenziato una regione di calibrazione ottimale compresa approssimativamente tra $T = 1.4$ e $T = 2.0$, in linea con quanto riportato nella letteratura sui modelli moderni.

Tuttavia, l’integrazione del classificatore all’interno della pipeline end-to-end di \textsc{PENTION-M} ha reso evidente come un valore di temperatura statico, seppur ottimale in media, non fosse sufficiente a garantire un comportamento affidabile in uno scenario operativo reale. Il sistema opera infatti in condizioni fortemente dinamiche, caratterizzate da variabilità atmosferica, turbolenza del vento, degradazione del segnale spettrale in funzione della distanza dalla sorgente, nonché da segnali diagnostici provenienti dalla pipeline MLOps, quali il drift statistico e la latenza operativa.

Alla luce di queste considerazioni, è stato progettato un meccanismo di \emph{Temperature Scaling adattivo}, concepito come una meta-euristica ibrida in grado di modulare il valore della temperatura in funzione di diverse fonti di incertezza. Il punto di partenza è rappresentato da una temperatura di base pari a $T=1.2$, alla quale vengono applicati incrementi controllati derivanti da:
\begin{itemize}
    \item il livello di rumore stimato nello spettro EI;
    \item la classe di stabilità atmosferica secondo Pasquill;
    \item indicatori di coerenza fisica provenienti dal modulo PIML di correzione della dispersione;
    \item il drift score calcolato dal servizio di monitoring MLOps;
    \item la latenza end-to-end della pipeline, utilizzata come proxy della congestione computazionale e della qualità del dato.
\end{itemize}

Il valore di temperatura risultante è vincolato a rimanere all’interno dell’intervallo $[1.0, 3.0]$ ed evolve in modo continuo e stabile nel tempo, evitando oscillazioni brusche e garantendo una progressione coerente rispetto alle condizioni operative. Durante le simulazioni integrate, la temperatura adattiva ha assunto valori tipicamente compresi tra $1.2$ e $2.6$, dimostrando un comportamento robusto e interpretabile.

L’introduzione del Temperature Scaling adattivo ha consentito di superare in modo definitivo il problema della sovrastima probabilistica, producendo confidenze più realistiche e sensibili al contesto fisico e computazionale della pipeline. Inoltre, il valore della temperatura utilizzata in ogni inferenza viene registrato dal Monitoring Service e incluso nel bundle forense generato dal sistema, garantendo la completa tracciabilità del processo di calibrazione. Ciò consente, in fase di audit, di ricostruire non solo la predizione finale, ma anche il livello di incertezza che ha accompagnato l’analisi, soddisfacendo i requisiti di trasparenza e verificabilità richiesti in ambito investigativo.

Nel complesso, questa fase ha trasformato il classificatore NPS da un modello affetto da overconfidence strutturale ad un componente pienamente integrato in una pipeline PIML+MLOps, dotato di un meccanismo di calibrazione probabilistica scientificamente fondato, validato sperimentalmente e coerente con le esigenze forensi di \textsc{PENTION-M}.
